% ----------------------
\subsection{Section 29 (September 1830)}
% ----------------------
	\label{DC29}
	
	% -----
	\subsubsection{Connection to the book of Moses}
	% -----
	
		Beginning around verse 29, there are some interesting connections between this section and the first few books of Moses, and in particular, \hyperref[MOSES-4]{Moses 4}.
	
	% -----
	\subsubsection{Historical Background}
	% -----
		\label{DC29-HB}
		
		% --
		\paragraph{JSP}
		% --
		
			``This revelation addressed the interest of some early church members in a Book of Mormon prophecy that described the physical gathering of God's chosen people in America. The Book of Mormon explained that during Christ's ministry in the Americas he prophesied that his chosen people would establish a sacred city, the New Jerusalem. According to the prophecies, ``the remnant of Jacob,'' which early church members identified as the American Indians, ``and also, as many of the house of Israel as shall come'' were to build this sacred city and gather to it, assisted by Gentiles who embraced the book's message. Christ further prophesied that when the progeny of the people described in the Book of Mormon were taught ``this Gospel'' again, Zion would be established among them.
			
			``According to the heading John Whitmer gave this text in Revelation Book 1, the setting for this revelation was a gathering of ``Six Elders of the Church \& three members'' who ``understood from Holy Writ that the time had come that the People of God should see eye to eye.'' The book of Isaiah declared that God's people would ``see eye to eye, when the Lord shall bring again Zion''; the Book of Mormon expressed the same sentiment and located Zion in the Americas. The heading seems to indicate, then, that this small group, believing that the Book of Mormon prophecy about Zion would soon be fulfilled, therefore ``enquired of the Lord \& thus came the word of the Lord through Joseph the seer.''
			
			``The revelation affirmed the imminent advent of the Millennium and declared that members of the Church of Christ were called to help gather God's people before the great event. It then turned to the creation of the world and the nature of Adam's fall, subjects JS had recently taken up in his Bible revision. According to the heading, the small group had differing views about ``the death of Adam (that is his transgression).'' Near the end of the text, the revelation addressed the question of whether God's commandment to Adam to not partake of the forbidden fruit was spiritual or temporal by declaring, ``All things unto me are Spiritual \& not at any time have I given unto you a law which was temporal neither any man nor the childern of men Neither Adam your father whom I created.'' Thus Adam's ``temporal'' act of eating the forbidden fruit rendered him ``spiritually dead.''
			
			``This revelation called for the gathering of God's people at the same time that a significant controversy had emerged among the membership of the Church of Christ. In September 1830, JS was attempting to address the problems arising from Hiram Page announcing his own revelations, the authenticity of which was accepted by a number of prominent church members, including Oliver Cowdery and the Whitmer family. Page's revelations, which concerned ``the upbuilding of Zion, the order of the Church \&c \&c,'' and this revelation's call to gather God's chosen people prompted another September revelation that clarified JS's prophetic role as the sole revelator for the church, required Cowdery to correct Hiram Page, and called Cowdery to preach to American Indians in the West.''%
				\footnote{%
					% \label{}%
					That's \hyperref[DC28]{DC 28}.
					}
			
	% -----
	\subsubsection{Versions}
	% -----
		\label{DC29-V}
		
		See the \href{https://drive.google.com/file/d/1goAvNz8jS4R8slYfMG_db55Ty2z_vmgK/view?usp=sharing}{sections comparison} document.
	
		\begin{compactdesc}
			\item[JSP] \hyperref[EMS-1830-JS-SEP-1830-A]{September 1830}
			\item[{\hyperref[HSN]{HSN}}\footnotemark] 29--38
			\item[EMS] 1:4 (September 1832), 26
			\item[BC] Chapter 29, 61--67
			\item[DC 1835] Section 10, 112--116
			\item[TS] 4:9 (15 March 1843), 130--131
			\item[MS] 4:11 (March 1844), 165--167
			\item[DC 1844] Section 10, 152--158
		\end{compactdesc}
			\footnotetext{%
				% \label{}%
				This source is the ``Hyde and Smith, Notebook.'' See \hyperref[HSN]{here}. According to the JSP website, it was copied down ``circa 25 January and 1 February 1832,'' by Orson Hyde. HDDC 417 says the date was 16 November 1832.
				}

	% -----
	\subsubsection{Section 29:1} % *DC29-1 
	% -----
		\label{DC29-1}

			\footnote{%
				% \label{}%
				In RB1, John Whitmer added this heading: ``A Revelation to Six Elders of the Church \& three members they understood from Holy Writ that the time had come <that> the People of God should see eye to eye \& they seeing somewhat different upon the death of Adam (that is his transgression) therefor they made it a subject of Prayer \& enquired of the Lord \& thus came the word of the Lord through Joseph the seer <saying given> At Fayette Seneca County State of New York.'' The JSP version is from RB1.
				}%
		LISTEN to the voice of Jesus Christ, your Redeemer, the Great I AM,%
			\footnote{%
				% \label{}%
				See \hyperref[EXODUS-3-14]{Exodus 3:14}. The phrase ``Great I AM'' appears three times in the scriptures, all in the DC---here, \hyperref[DC38-1]{38:1}, and \hyperref[DC39-1]{39:1}.
				}
		whose arm of mercy%
			\footnote{%
				% \label{}%
				``Arm of mercy'' is a restoration construction, occuring seven total times. See \hyperref[MERCY-PHRASES]{here} for a list.
				}
		hath atoned for your sins;

		% \begin{framed}
		% \scriptsize
			% \begin{compactdesc}
				% \item[JSP] 
				% % \item[HSN] 
				% % \item[EMS] 
				% % \item[BC] 
				% % \item[]
			% \end{compactdesc}
		% \end{framed}

	% -----
	\subsubsection{Section 29:2} % *DC29-2 
	% -----
		\label{DC29-2}

		Who will gather his people even as a hen gathereth her chickens under her wings,%
			\footnote{%
				% \label{}%
				See \hyperref[MATTHEW-23-37]{Matthew 23:37}.
				}
		even as many as will hearken to my voice and humble themselves before me, and call upon me in mighty prayer.

		% \begin{framed}
		% \scriptsize
			% \begin{compactdesc}
				% \item[JSP] 
				% % \item[HSN] 
				% % \item[EMS] 
				% % \item[BC] 
				% % \item[]
			% \end{compactdesc}
		% \end{framed}

	% -----
	\subsubsection{Section 29:3} % *DC29-3 
	% -----
		\label{DC29-3}

		Behold, verily, verily, I say unto you, that at this time your sins are forgiven you, therefore ye receive these things; but remember to sin no more, lest perils shall come upon you.

		\begin{framed}
		\scriptsize
			\begin{compactdesc}
				\item[JSP] Behold Verily Verily I say unto you at this time your sins are forgiven you Therefore ye Receive these things but remember to sin no more lest perils shall come upon you
				\item[HSN] behold verily I say unto you that at this time your sins are forgiven you therefore ye receive these things but remember to sin no more lest perils shall come upon you
				% \item[EMS] 
				% \item[BC] 
				% \item[]
			\end{compactdesc}
		\end{framed}

	% -----
	\subsubsection{Section 29:4} % *DC29-4 
	% -----
		\label{DC29-4}

		Verily, I say unto you that ye are chosen out of the world to declare my gospel with the sound of rejoicing, as with the voice of a trump.

		% \begin{framed}
		% \scriptsize
			% \begin{compactdesc}
				% \item[JSP] 
				% % \item[HSN] 
				% % \item[EMS] 
				% % \item[BC] 
				% % \item[]
			% \end{compactdesc}
		% \end{framed}

	% -----
	\subsubsection{Section 29:5} % *DC29-5 
	% -----
		\label{DC29-5}

		Lift up your hearts and be glad, for I am in your midst, and am your \hyperref[ADVOCATE]{advocate} with the Father; and it is his good will to give you the kingdom.%
			\footnote{%
				% \label{}%
				See \hyperref[DC76-5]{DC 76:5}.
				}

		% \begin{framed}
		% \scriptsize
			% \begin{compactdesc}
				% \item[JSP] 
				% % \item[HSN] 
				% % \item[EMS] 
				% % \item[BC] 
				% % \item[]
			% \end{compactdesc}
		% \end{framed}

	% -----
	\subsubsection{Section 29:6} % *DC29-6 
	% -----
		\label{DC29-6}

		And, as it is written---Whatsoever ye shall ask in faith, being united in prayer according to my command, ye shall receive.%
			\footnote{%
				% \label{}%
				There are various versions of this promise in the scriptures---the ones most similar to what is said here, since it says ``as it is written,'' are at \hyperref[ENOS-1-15]{Enos 1:15}, \hyperref[MOSIAH-4-21]{Mosiah 4:21}, \hyperref[MORONI-7-26]{Moroni 7:26}, \hyperref[DC8-1]{DC 8:1}, and \hyperref[DC27-18]{DC 27:18}.
				}

		% \begin{framed}
		% \scriptsize
			% \begin{compactdesc}
				% \item[JSP] 
				% % \item[HSN] 
				% % \item[EMS] 
				% % \item[BC] 
				% % \item[]
			% \end{compactdesc}
		% \end{framed}

	% -----
	\subsubsection{Section 29:7} % *DC29-7 
	% -----
		\label{DC29-7}

		And ye are called to bring to pass the gathering of mine \hyperref[ELECTION]{elect}; for mine \hyperref[ELECTION]{elect} hear my voice and harden not their hearts;

		% \begin{framed}
		% \scriptsize
			% \begin{compactdesc}
				% \item[JSP] 
				% % \item[HSN] 
				% % \item[EMS] 
				% % \item[BC] 
				% % \item[]
			% \end{compactdesc}
		% \end{framed}

	% -----
	\subsubsection{Section 29:8} % *DC29-8 
	% -----
		\label{DC29-8}

		Wherefore the decree hath gone forth from the Father that they shall be gathered in unto one place upon the face of this land, to prepare their hearts and be prepared in all things against the day when tribulation and desolation are sent forth upon the wicked.

		% \begin{framed}
		% \scriptsize
			% \begin{compactdesc}
				% \item[JSP] 
				% % \item[HSN] 
				% % \item[EMS] 
				% % \item[BC] 
				% % \item[]
			% \end{compactdesc}
		% \end{framed}

	% -----
	\subsubsection{Section 29:9} % *DC29-9 
	% -----
		\label{DC29-9}

		For the hour is nigh and the day soon at hand when the earth is ripe; and all the proud and they that do wickedly shall be as stubble; and I will burn them up, saith the Lord of Hosts, that wickedness shall not be upon the earth;

		% \begin{framed}
		% \scriptsize
			% \begin{compactdesc}
				% \item[JSP] 
				% % \item[HSN] 
				% % \item[EMS] 
				% % \item[BC] 
				% % \item[]
			% \end{compactdesc}
		% \end{framed}

	% -----
	\subsubsection{Section 29:10} % *DC29-10 
	% -----
		\label{DC29-10}

		For the hour is nigh, and that which was spoken by mine apostles must be fulfilled; for as they spoke so shall it come to pass;

		% \begin{framed}
		% \scriptsize
			% \begin{compactdesc}
				% \item[JSP] 
				% % \item[HSN] 
				% % \item[EMS] 
				% % \item[BC] 
				% % \item[]
			% \end{compactdesc}
		% \end{framed}

	% -----
	\subsubsection{Section 29:11} % *DC29-11 
	% -----
		\label{DC29-11}

		For I will reveal myself from heaven with power and great glory, with all the hosts thereof, and dwell in righteousness with men on earth a thousand years, and the wicked shall not stand.

		% \begin{framed}
		% \scriptsize
			% \begin{compactdesc}
				% \item[JSP] 
				% % \item[HSN] 
				% % \item[EMS] 
				% % \item[BC] 
				% % \item[]
			% \end{compactdesc}
		% \end{framed}

	% -----
	\subsubsection{Section 29:12} % *DC29-12 
	% -----
		\label{DC29-12}

		And again, verily, verily, I say unto you, and it hath gone forth in a firm decree, by the will of the Father, that mine apostles, the Twelve which were with me in my ministry at Jerusalem, shall stand at my right hand at the day of my coming in a pillar of fire, being clothed with robes of righteousness, with crowns upon their heads,%
			\footnote{%
				% \label{}%
				Sounds like temple clothing. See \hyperref[JOHN-19-2]{John 19:2}, and \hyperref[DC109-76]{DC 109:76}.
				}
		in glory even as I am,%
			\footnote{%
				% \label{}%3
				The apostles have received glory like Christ's. See \hyperref[DC93-22]{DC 93:22, 28}.
				}
		to judge the whole house of Israel, even as many as have loved me and kept my commandments, and none else.

		\begin{framed}
		\scriptsize
			\begin{compactdesc}
				\item[JSP] \& again Verily Verily I say unto you \& it hath gone forth in a firm decree by the will of the father that mine Apostles the twelve which were with me in my ministery at Jerusalem shall stand at my right hand at the day of my comeing in a piller of fire being clothed with robes of righteousness with crowns upon their heads in glory even as I am to Judge the whole House of Israel even as many as have loved me \& kept my commandments \& none else
				\item[HSN] and again verily I say unto you \& it hath gone forth in a firm decree by the will of the father that mine apostles the 12 which were with me in my ministry at Jerusalem shall stand at my right hand at the day of my coming in a pillar of fire being clothed with robes of righteousness with \sout{crowith} crowns upon their heads in glory even as I am to judge the whole house of Israel even as many as have loved me and kept my commandments and none else
				% \item[EMS] 
				% \item[BC] 
				% \item[]
			\end{compactdesc}
		\end{framed}

	% -----
	\subsubsection{Section 29:13} % *DC29-13 
	% -----
		\label{DC29-13}

		For a trump shall sound both long and loud, even as upon Mount Sinai, and all the earth shall quake, and they shall come forth---yea, even the dead which died in me, to receive a crown of righteousness, and to be clothed upon, even as I am, to be with me, that we may be one.

		% \begin{framed}
		% \scriptsize
			% \begin{compactdesc}
				% \item[JSP] 
				% % \item[HSN] 
				% % \item[EMS] 
				% % \item[BC] 
				% % \item[]
			% \end{compactdesc}
		% \end{framed}

	% -----
	\subsubsection{Section 29:14} % *DC29-14 
	% -----
		\label{DC29-14}

		But, behold, I say unto you that before this great day shall come the sun shall be darkened, and the moon shall be turned into blood, and the stars shall fall from heaven, and there shall be greater signs in heaven above and in the earth beneath;

		% \begin{framed}
		% \scriptsize
			% \begin{compactdesc}
				% \item[JSP] 
				% % \item[HSN] 
				% % \item[EMS] 
				% % \item[BC] 
				% % \item[]
			% \end{compactdesc}
		% \end{framed}

	% -----
	\subsubsection{Section 29:15} % *DC29-15 
	% -----
		\label{DC29-15}

		And there shall be weeping and wailing among the hosts of men;

		% \begin{framed}
		% \scriptsize
			% \begin{compactdesc}
				% \item[JSP] 
				% % \item[HSN] 
				% % \item[EMS] 
				% % \item[BC] 
				% % \item[]
			% \end{compactdesc}
		% \end{framed}

	% -----
	\subsubsection{Section 29:16} % *DC29-16 
	% -----
		\label{DC29-16}

		And there shall be a great hailstorm sent forth to destroy the crops of the earth.

		% \begin{framed}
		% \scriptsize
			% \begin{compactdesc}
				% \item[JSP] 
				% % \item[HSN] 
				% % \item[EMS] 
				% % \item[BC] 
				% % \item[]
			% \end{compactdesc}
		% \end{framed}

	% -----
	\subsubsection{Section 29:17} % *DC29-17 
	% -----
		\label{DC29-17}

		And it shall come to pass, because of the wickedness of the world, that I will take vengeance upon the wicked, for they will not repent; for the cup of mine indignation is full; for behold, my blood shall not cleanse them if they hear me not.%
			\footnote{%
				% \label{}%
				See, for example, \hyperref[ALMA-13-11]{Alma 13:11}.
				}

		% \begin{framed}
		% \scriptsize
			% \begin{compactdesc}
				% \item[JSP] 
				% % \item[HSN] 
				% % \item[EMS] 
				% % \item[BC] 
				% % \item[]
			% \end{compactdesc}
		% \end{framed}

	% -----
	\subsubsection{Section 29:18} % *DC29-18 
	% -----
		\label{DC29-18}

		Wherefore, I the Lord God will send forth flies upon the face of the earth, which shall take hold of the inhabitants thereof, and shall eat their flesh, and shall cause maggots to come in upon them;

		% \begin{framed}
		% \scriptsize
			% \begin{compactdesc}
				% \item[JSP] 
				% % \item[HSN] 
				% % \item[EMS] 
				% % \item[BC] 
				% % \item[]
			% \end{compactdesc}
		% \end{framed}

	% -----
	\subsubsection{Section 29:19} % *DC29-19 
	% -----
		\label{DC29-19}

		And their tongues shall be stayed that they shall not utter against me; and their flesh shall fall from off their bones, and their eyes from their sockets;

		% \begin{framed}
		% \scriptsize
			% \begin{compactdesc}
				% \item[JSP] 
				% % \item[HSN] 
				% % \item[EMS] 
				% % \item[BC] 
				% % \item[]
			% \end{compactdesc}
		% \end{framed}

	% -----
	\subsubsection{Section 29:20} % *DC29-20 
	% -----
		\label{DC29-20}

		And it shall come to pass that the beasts of the forest and the fowls of the air shall devour them up.

		% \begin{framed}
		% \scriptsize
			% \begin{compactdesc}
				% \item[JSP] 
				% % \item[HSN] 
				% % \item[EMS] 
				% % \item[BC] 
				% % \item[]
			% \end{compactdesc}
		% \end{framed}

	% -----
	\subsubsection{Section 29:21} % *DC29-21 
	% -----
		\label{DC29-21}

		And the great and abominable church, which is the whore of all the earth, shall be cast down by devouring fire, according as it is spoken by the mouth of Ezekiel the prophet, who spoke of these things,%
			\footnote{%
				% \label{}%
				Possibly \hyperref[EZEKIEL-24-6]{Ezekiel 24:6--14}?
				}
		which have not come to pass but surely must, as I live, for abominations shall not reign.

		% \begin{framed}
		% \scriptsize
			% \begin{compactdesc}
				% \item[JSP] 
				% % \item[HSN] 
				% % \item[EMS] 
				% % \item[BC] 
				% % \item[]
			% \end{compactdesc}
		% \end{framed}

	% -----
	\subsubsection{Section 29:22} % *DC29-22 
	% -----
		\label{DC29-22}

		And again, verily, verily, I say unto you that when the thousand years are ended, and men again begin to deny their God, then will I spare the earth but for a little season;

		% \begin{framed}
		% \scriptsize
			% \begin{compactdesc}
				% \item[JSP] 
				% % \item[HSN] 
				% % \item[EMS] 
				% % \item[BC] 
				% % \item[]
			% \end{compactdesc}
		% \end{framed}

	% -----
	\subsubsection{Section 29:23} % *DC29-23 
	% -----
		\label{DC29-23}

		And the end shall come, and the heaven and the earth shall be consumed and pass away, and there shall be a new heaven and a new earth.

		% \begin{framed}
		% \scriptsize
			% \begin{compactdesc}
				% \item[JSP] 
				% % \item[HSN] 
				% % \item[EMS] 
				% % \item[BC] 
				% % \item[]
			% \end{compactdesc}
		% \end{framed}

	% -----
	\subsubsection{Section 29:24} % *DC29-24 
	% -----
		\label{DC29-24}

		For all old things shall pass away, and all things shall become new, even the heaven and the earth, and all the fulness thereof, both men and beasts, the fowls of the air, and the fishes of the sea;

		% \begin{framed}
		% \scriptsize
			% \begin{compactdesc}
				% \item[JSP] 
				% % \item[HSN] 
				% % \item[EMS] 
				% % \item[BC] 
				% % \item[]
			% \end{compactdesc}
		% \end{framed}

	% -----
	\subsubsection{Section 29:25} % *DC29-25 
	% -----
		\label{DC29-25}

		And not one hair, neither mote, shall be lost, for it is the workmanship of mine hand.

		% \begin{framed}
		% \scriptsize
			% \begin{compactdesc}
				% \item[JSP] 
				% % \item[HSN] 
				% % \item[EMS] 
				% % \item[BC] 
				% % \item[]
			% \end{compactdesc}
		% \end{framed}

	% -----
	\subsubsection{Section 29:26} % *DC29-26 
	% -----
		\label{DC29-26}

		But, behold, verily I say unto you, before the earth shall pass away, Michael, mine archangel, shall sound his trump, and then shall all the dead awake, for their graves shall be opened, and they shall come forth---yea, even all.

		% \begin{framed}
		% \scriptsize
			% \begin{compactdesc}
				% \item[JSP] 
				% % \item[HSN] 
				% % \item[EMS] 
				% % \item[BC] 
				% % \item[]
			% \end{compactdesc}
		% \end{framed}

	% -----
	\subsubsection{Section 29:27} % *DC29-27 
	% -----
		\label{DC29-27}

		And the righteous shall be gathered on my right hand unto eternal life; and the wicked on my left hand will I be ashamed to own before the Father;

		% \begin{framed}
		% \scriptsize
			% \begin{compactdesc}
				% \item[JSP] 
				% % \item[HSN] 
				% % \item[EMS] 
				% % \item[BC] 
				% % \item[]
			% \end{compactdesc}
		% \end{framed}

	% -----
	\subsubsection{Section 29:28} % *DC29-28 
	% -----
		\label{DC29-28}

		Wherefore I will say unto them---Depart from me, ye cursed, into everlasting fire, prepared for the devil and his angels.

		% \begin{framed}
		% \scriptsize
			% \begin{compactdesc}
				% \item[JSP] 
				% % \item[HSN] 
				% % \item[EMS] 
				% % \item[BC] 
				% % \item[]
			% \end{compactdesc}
		% \end{framed}

	% -----
	\subsubsection{Section 29:29} % *DC29-29 
	% -----
		\label{DC29-29}

		And now, behold, I say unto you, never at any time have I declared from mine own mouth that they should return, for where I am they cannot come, for they have no power.%
			\footnote{%
				% \label{}%
				What ``power'' is the Lord speaking of here? Some kind of priesthood power. Could just be referring to the \hyperref[PRIESTHOOD-powers]{power of resurrection}, or maybe the \hyperref[PRIESTHOOD-powers]{power of atonement}, since in \hyperref[2NEPHI-10-25]{2 Nephi 10:25} that power is spoken of as the power that will allow a person to ``be received into the eternal kingdom of God.''
				}

		\begin{framed}
		\scriptsize
			\begin{compactdesc}
				\item[JSP] \& now Behold I say unto you never at any time have I declared from mine own mouth that they should return for where I am they cannot come for they have no power
				\item[HSN] and now behold I say unto you never at any time hath I declard from mine own mouth that they should return for where I am they cannot come for they have no power,
				% \item[EMS] 
				% \item[BC] 
				% \item[]
			\end{compactdesc}
		\end{framed}

	% -----
	\subsubsection{Section 29:30} % *DC29-30 
	% -----
		\label{DC29-30}

		But remember that all my judgments are not given unto men; and as the words have gone forth out of my mouth even so shall they be fulfilled, that the first shall be last, and that the last shall be first in all things whatsoever I have created by the word of my power,%
		\footnote{%
			% \label{}%
			For occurrences of ``word of my power'' or very similar phrases, see \hyperref[WORD-PHRASES]{here}.
			} 
		which is the power of my Spirit.%
			\footnote{%
				% \label{}%
				For more occurrences of phrases similar to ``power of my Spirit,'' see \hyperref[POWER-PHRASES]{here}. Note the very important connection to \hyperref[DC88-7]{DC 88:7--13}. Those verses do not use the word ``spirit,'' but they talk about the ``light of Christ'' (verse 7), which is then described as ``the power of God who sitteth upon his throne.'' I think the key is recognizing that the ``light of Christ'' is the same as the ``spirit of Christ,'' and that that ``spirit'' is what this verse in DC 29:30 is speaking of. See \hyperref[LIGHT-christ]{this study} for more.
				} 

		\begin{framed}
		\scriptsize
			\begin{compactdesc}
				\item[JSP] but remember that all my Judgements are not given unto men \& as the words have gone forth out of my mouth even so shall they be fulfilled that the first shall be last \& that the last shall be first in all things Whatsoever I have created by the word of my Power which is the Power of my spirit
				\item[HSN] but remember that all my judgments are not given unto man, and as the words have gone forth out of my mouth even so shall they be fulfilled that the first should be last and that the last should be first in all things, whatsoever I have created by the word of my power which is the power of my spirit
				% \item[EMS] 
				% \item[BC] 
				% \item[]
			\end{compactdesc}
		\end{framed}

	% -----
	\subsubsection{Section 29:31} % *DC29-31 
	% -----
		\label{DC29-31}

		For by the power of my Spirit created I them; yea, all things both spiritual and temporal---%
			\footnote{%
				% \label{}%
				The spiritual/temporal distinction which appears here and in the next few verses reminds me of the ``spiritual creation'' described in \hyperref[MOSES-3-4]{Moses 3:4--7}.
				}

		\begin{framed}
		\scriptsize
			\begin{compactdesc}
				\item[JSP] for by the Power of my Spirit created I them yea all things both Spiritual \& Temperally
				\item[HSN] for by the power of my spirit created I them, yea all things both spiritual and temporal
				% \item[EMS] 
				% \item[BC] 
				% \item[]
			\end{compactdesc}
		\end{framed}

	% -----
	\subsubsection{Section 29:32} % *DC29-32 
	% -----
		\label{DC29-32}

		First \hyperref[SPIRITUAL]{spiritual}, secondly \hyperref[TEMPORAL]{temporal}, which is the beginning of my work; and again, first \hyperref[TEMPORAL]{temporal}, and secondly \hyperref[SPIRITUAL]{spiritual}, which is the last of my work---

		\begin{framed}
		\scriptsize
			\begin{compactdesc}
				\item[JSP] firstly spiritual secondly temporally which is the Begining of my work \& again firstly temporal \& secondly spiritual which is the last of my work
				\item[HSN] firstly spiritual secondly temporal which is the beginning of my work, and again firstly temporal and secondly spiritual which is the last of my work
				% \item[EMS] 
				% \item[BC] 
				% \item[]
			\end{compactdesc}
		\end{framed}

	% -----
	\subsubsection{Section 29:33} % *DC29-33 
	% -----
		\label{DC29-33}

		Speaking unto you that you may naturally understand; but unto myself my works have no end, neither beginning; but it is given unto you that ye may understand, because ye have asked it of me and are agreed.

		\begin{framed}
		\scriptsize
			\begin{compactdesc}
				\item[JSP] speaking unto you that ye may naturally understand but unto myself my work hath no end neither begining But it is given unto you that ye may understand because ye have asked it of me \& are agreed
				\item[HSN] speaking unto you that ye may naturally understand, but unto myself my works hath no end neither beginning but it is given unto you that ye may understand because ye have asked it of me and are agreed
				% \item[EMS] 
				% \item[BC] 
				% \item[]
			\end{compactdesc}
		\end{framed}

	% -----
	\subsubsection{Section 29:34} % *DC29-34 
	% -----
		\label{DC29-34}

		Wherefore, verily I say unto you that all things unto me are \hyperref[SPIRITUAL]{spiritual}, and not at any time have I given unto you a \hyperref[LAW]{law} which was \hyperref[TEMPORAL]{temporal};%
			\footnote{%
				% \label{}%
				See \hyperref[2NEPHI-2-5]{2 Nephi 2:5} and notes there.
				}
		neither any man, nor the children of men; neither Adam, your father, whom I created.

		\begin{framed}
		\scriptsize
			\begin{compactdesc}
				\item[JSP] Wherefore Verily I say unto you that all things unto me are Spiritual \& not at any time have I given unto you a law which was temporal neither any man nor the childern of men Neither Adam your father whom I created
				\item[HSN] wherefore verily I say unto you that all things unto me are spiritual and not at any time have I given unto you a law which was temporal neither any man neither the children of men neither adam your father whom I have created
				% \item[EMS] 
				% \item[BC] 
				% \item[]
			\end{compactdesc}
		\end{framed}

	% -----
	\subsubsection{Section 29:35} % *DC29-35 
	% -----
		\label{DC29-35}

		Behold, I gave unto him%
			\footnote{%
				% \label{}%
				I take the ``I gave unto him'' part to be something like, ``I prepared the conditions in which he could be an agent unto himself.''
				}
		that he should be an \hyperref[AGENCY]{agent} unto himself;%
			\footnote{%
				% \label{}%
				This whole first part of the verse seems equivalent to the first part of \hyperref[2NEPHI-2-16]{2 Nephi 2:16}: ``Wherefore the Lord God gave unto man that he should act for himself.''
				}
		and I gave unto him commandment,%
			\footnote{%
				% \label{}%
				Note that the JSP version has an indefinite article here, ``\textit{a} commandment.''
				}
		but no \hyperref[TEMPORAL]{temporal} commandment gave I unto him, for my commandments are \hyperref[SPIRITUAL]{spiritual}; they are not \hyperref[NATURAL]{natural} nor \hyperref[TEMPORAL]{temporal}, neither \hyperref[CARNAL]{carnal} nor \hyperref[SENSUAL]{sensual}.%
			\footnote{%
				% \label{}%
				See \hyperref[JAMES-3-15]{James 3:15} and notes there.
				}

		\begin{framed}
		\scriptsize
			\begin{compactdesc}
				\item[JSP] Behold I \sout{give} <gave> unto him that he should be an agent unto himself \& I gave unto him a commandment but no temporal Commandment gave I unto him for my commandments are spiritual they are not Natural nor temporal neither carnal nor sensual
				\item[HSN] behold I gave unto him but no carnal commandments for my commandments are spiritual th[e]y are not natural nor temporal neither carnal nor sensual
				% \item[EMS] 
				% \item[BC] 
				% \item[]
			\end{compactdesc}
		\end{framed}

	% -----
	\subsubsection{Section 29:36} % *DC29-36 
	% -----
		\label{DC29-36}

		And it came to pass that Adam, being tempted of the devil---for, behold, the devil was before Adam,%
			\footnote{%
				% \label{}%
				Existed before, or was standing in front of, or both (or some other meaning).
				}
		for he rebelled against me,%
		\footnote{%
			% \label{}%
			See \hyperref[DC76-25]{DC 76:25} and \hyperref[DC76-28]{28}, but especially \hyperref[MOSES-4-3]{Moses 4:3}, which is a very important parallel to this verse.
			} 
		saying, Give me thine \hyperref[HONOR]{honor}, which is my power;%
			\footnote{%
				% \label{}%
				I think the ``honor'' spoken of here is God's priesthood. Compare \hyperref[HEBREWS-5-4]{Hebrews 5:4}, where you get \hyperref[TIME-2]{\grl{τιμή}}; see definition \#2.d in BDAG for that word, and also meanings \#2 and \#3 in CGL. The relative clause which immediately follows the statement tells us that the ``honor'' is God's power, and his power is the priesthood. Some people, like Cleon Skousen in an essay somewhere, understand this passage as saying that the honor God receives is his power, but I think that is wrong. The word ``honor'' is just calling attention to a particular aspect of God's power, but I think the term still refers here to the priesthood; compare \hyperref[MOSES-4-1]{Moses 4:1} and \hyperref[MOSES-4-3]{4:3}. Satan wanted the privilege, honor, and glory that God possesses in virtue of having that priesthood power.
				}
		and also a third part%
			\footnote{%
				% \label{}%
				I heard Hyrum Andrus once say in a lecture that the phrase ``a third part'' refers not to a fraction of the total amount, $\frac{1}{3}$, but to the third part out of three different parts, so ``third'' isn't a fraction but an ordinal number. This is the only occurrence of ``third part'' in the restoration scriptures; the phrase occurs 35 times in the Bible, including 14 times in Revelation (the only NT occurrences), and 21 times in the OT. \hyperref[REVELATIONNT-8-7]{Revelation 8:7--12} uses the phrase 11 times; it occurs once each at \hyperref[REVELATIONNT-9-15]{Revelation 9:15} and \hyperref[REVELATIONNT-9-18]{Revelation 9:18}; and then at \hyperref[REVELATIONNT-12-4]{Revelation 12:4}. The most relevant parallel to this verse in DC 29 is Revelation 12:4 I think. The Greek for ``third part'' is \gr{τὸ τρίτον} (that phrase is also all over Revelation 8). \gr{τρίτος} is at BDAG 903, where the first meaning is ``\textbf{third in a series, \textit{third}},'' but this is when the word is used as an adjective, which is not the Revelation uses. Used substantively as a noun phrase, as it is in Revelation, BDAG meaning \#2 says ``as a substantive, \textbf{a third part of something, \textit{third, third part}}.'' Later in the entry it clarifies that this is the fraction use: ``\textit{the third part, one-third} followed by partitive genitive,'' after which it lists all the Revelation uses I mentioned above, including 12:4. Note that you get a partitive genitive in 12:4: \gr{τὸ τρίτον τῶν ἀστέρων}. So the key seems to be whether the word is used as an adjective, frequently in attributive position (\hyperref[REVELATIONNT-8-10]{Revelation 8:10}, \gr{ὁ τρίτος ἄγγελος}), where it means ``third in a series,'' (the ordinal meaning), or is used as a substantive, where it's a way of saying the fraction ``one-third.'' Compare the alternate translations at \hyperref[REVELATIONNT-12-4]{12:4}. Given this evidence, the strongest claim you can make against Andrus is that, if DC 29:36 is intended to be a parallel to Revelation 12:4, then it can't mean ``the third part of a group of three items.'' I don't see any reason to think that this is a restoration innovation; it would be simpler to think that roughly $\frac{1}{3}$ of the individuals in the overall group were drawn away. (Compare CGL 1396a, meaning \#1, ``(of persons or things) \textbf{third} (in a list or sequence, often as last item),'' with meaning \#5, ``(phrase) \gr{τρίτον μέρος} \textit{third part, one third} (\textit{of an object, country, group of persons, or similar}).'') 
				}
		of the hosts of heaven turned he away from me because of their \hyperref[AGENCY]{agency};%
			\footnote{%
				% \label{}%
				See \hyperref[MOSES-4-1]{Moses 4:1}. Satan, apparently, wanted to obligate people to behave in a certain way such that they would be saved, or at least, they would do all the things on the outside in order to be saved.
				}

		\begin{framed}
		\scriptsize
			\begin{compactdesc}
				\item[JSP] \& it came to pass that Adam being tempted of the Devil for Behold the Devil was before Adam for he rebelled against me saying give me thine honour which is my Power \& also a third part of the host of Heaven turned he away from me Because of their agency
				\item[HSN] and it came to pass that adam being <tempted> of the Devil for behold the Devil was before Adam for he rebelled against me saying give me thine honour which is my power and also the Third part of the hosts of heaven turned he away from me because of their agency
				\item[EMS] and it came to pass, that Adam, being tempted of the Devil, for behold the Devil was before Adam, for he rebelled against me, saying, Give me thine honor, which is my power, and also a third part af [sic] the host of heaven turned he away from me because of their agency:
				\item[BC] 45 And it came to pass, that Adam being tempted of the devil, for behold the devil was before Adam, for he rebelled against me saying, Give me thine honor, which is my power: and also a third part of the hosts of heaven turned he away from me because of their agency:
				% \item[]
			\end{compactdesc}
		\end{framed}

	% -----
	\subsubsection{Section 29:37} % *DC29-37 
	% -----
		\label{DC29-37}

		And they were thrust down, and thus came the devil and his angels;

		\begin{framed}
		\scriptsize
			\begin{compactdesc}
				\item[JSP] \& they were thrust down \& thus came the Devil \& his Angels
				\item[HSN] and They were Thrust down \& thus came the Devil and his angels
				% \item[EMS] 
				% \item[BC] 
				% \item[]
			\end{compactdesc}
		\end{framed}

	% -----
	\subsubsection{Section 29:38} % *DC29-38 
	% -----
		\label{DC29-38}

		And, behold, there is a place prepared for them from the beginning, which place is hell.

		\begin{framed}
		\scriptsize
			\begin{compactdesc}
				\item[JSP] \& Behold a place prepared for them which place is Hell
				\item[HSN] \& behold a place prepared for them from the beginning which place is hell
				% \item[EMS] 
				% \item[BC] 
				% \item[]
			\end{compactdesc}
		\end{framed}

	% -----
	\subsubsection{Section 29:39} % *DC29-39 
	% -----
		\label{DC29-39}

		And it must needs be that the devil should \hyperref[TEMPTATION]{tempt} the children of men, or they could not be \hyperref[AGENCY]{agents} unto themselves;%
		\footnote{%
			% \label{}%
			Given what it says here and after this, there's a knowledge component to agency, among whatever other requirements there might be.
			} 
		for if they never should have bitter they could not know the sweet---%
			\footnote{%
				% \label{}%
				It doesn't say that people couldn't \textit{experience} the sweet; of course they could. What it means is that, without having the bitter, they couldn't know the sweet \textit{as} the sweet, or couldn't know \textit{that} what they were experiencing was something to be sought after and valued, because of how different it is from other things. See \hyperref[MOSES-6-55]{Moses 6:55}, which is clearer: ``they taste the bitter, that they may know to prize the good.'' See also \hyperref[2NEPHI-2-16]{2 Nephi 2:16}.
				}

		\begin{framed}
		\scriptsize
			\begin{compactdesc}
				\item[JSP] \& it \sout{came to pass} Must needs be that the Devil should tempt the children of men or they could not be agents unto themselves for if they never should have bitter they could not k[n]ow the Sweet
				\item[HSN] and must be that the Devil should tempt the children of men or they cannot be agents unto themselves for if they never should have bitter they would not know the sweet
				% \item[EMS] 
				% \item[BC] 
				% \item[]
			\end{compactdesc}
		\end{framed}

	% -----
	\subsubsection{Section 29:40} % *DC29-40 
	% -----
		\label{DC29-40}

		Wherefore, it came to pass that the devil \hyperref[TEMPTATION]{tempted} Adam, and he partook of the forbidden fruit and transgressed the commandment, wherein he became subject to the will of the devil,%
			\footnote{%
				% \label{}%
				See \hyperref[2NEPHI-2-29]{2 Nephi 2:29}.
				}
		because he yielded unto \hyperref[TEMPTATION]{temptation}.

		\begin{framed}
		\scriptsize
			\begin{compactdesc}
				\item[JSP] Wherefore it came to pass that the Devil tempted Adam \& he partook of the forbiden fruit \& transgressed the commandment wherein he became subject to the will of the Devil Because he yielded unto temptation
				\item[HSN] wherefore it came to pass that the Devil tempted Adam \& he partook of the forbidden fruit and transgressed the commandments wherefore he became subject to the will of the Devil because he yielded to temptation
				% \item[EMS] 
				% \item[BC] 
				% \item[]
			\end{compactdesc}
		\end{framed}

	% -----
	\subsubsection{Section 29:41} % *DC29-41 
	% -----
		\label{DC29-41}

		Wherefore, I, the Lord God, caused that he should be cast out from the \hyperref[GARDENOFEDEN]{Garden of Eden}, from my presence, because of his transgression, wherein he became spiritually dead,%
			\footnote{%
				% \label{}%
				He no longer had access to the life-giving or quickening powers that flow so freely from God's presence; see \hyperref[JOHN-4-14]{John 4:14} and \hyperref[DC88-11]{DC 88:11--13}.
				}
		which is the first death, even that same death which is the last death, which is spiritual, which shall be pronounced upon the wicked when I shall say: Depart, ye cursed.

		\begin{framed}
		\scriptsize
			\begin{compactdesc}
				\item[JSP] Wherefore I the Lord God caused that he should be cast out from the Garden of Edan from my presence because of his transgression Wherein he became spiritually dead which \sout{death} is the first death even that same death which is the last death which is spiritual \sout{which shall be} pronounced upon the wicked which shall be when I shall say depart ye Cursed
				\item[HSN] wherefore I the Lord God caused that he should be cast out from the garden of Eden from my presence because of his transgression wherein he became spiritally dead which is the first Death even that that same death which is the last death which is spiritual which shall be pronounced upon the wicked when I shall say depart ye cursed
				% \item[EMS] 
				% \item[BC] 
				% \item[]
			\end{compactdesc}
		\end{framed}

	% -----
	\subsubsection{Section 29:42} % *DC29-42 
	% -----
		\label{DC29-42}

		But, behold, I say unto you that I, the Lord God, gave unto Adam and unto his seed, that they should not die as to the temporal death, until I, the Lord God, should send forth \hyperref[ANGEL]{angels} to declare unto them repentance and redemption, through faith on the name of mine Only Begotten Son.

		\begin{framed}
		\scriptsize
			\begin{compactdesc}
				\item[JSP] But Behold I say unto you that I the Lord God gave unto Adam \& unto his seed that they should not Die as to the temporal death untill I the Lord God should send forth Angels to declare unto them Repentance \& redemption through faith on the name of mine only begotten Son
				\item[HSN] but behold I say unto you that I the Lord God gave unto Adam \& unto his seed that they should not die as to the temporal death until I the Lord God should send forth Angels to declare unto them repentance \& redemption through faith on the name of mine onely begotten son
				% \item[EMS] 
				% \item[BC] 
				% \item[]
			\end{compactdesc}
		\end{framed}

	% -----
	\subsubsection{Section 29:43} % *DC29-43 
	% -----
		\label{DC29-43}

		And thus did I, the Lord God, appoint%
			\footnote{%
				% \label{}%
				Other similar uses of ``appoint'' include \hyperref[PSALMS-102-20]{Psalms 102:20}, \hyperref[1CORINTHIANS-4-9]{1 Corinthians 4:9}, \hyperref[ALMA-12-27]{Alma 12:27}, \hyperref[ALMA-40-7]{Alma 40:7}, \hyperref[ALMA-40-9]{Alma 40:9}, \hyperref[ALMA-40-21]{Alma 40:21}, and \hyperref[DC42-48]{DC 42:48}. The Greek word in the Corinthians verse is a form of \gr{ἀποδείκνυμι}.
				}
		unto man the days of his \hyperref[PROBATION]{probation}%
			\footnote{%
				% \label{}%
				For more occurrences of ``day(s) of probation'' \hyperref[PROBATION-PHRASES]{here}.
				}%
		---that by his \hyperref[NATURAL]{natural} death%
			\footnote{%
				% \label{}%
				This use of ``natural'' reminds me of \hyperref[MOSIAH-3-16]{Mosiah 3:16}; compare the next verse here, the ``spiritual fall.''
				}
		he might be raised in immortality unto eternal life, even as many as would believe;

		\begin{framed}
		\scriptsize
			\begin{compactdesc}
				\item[JSP] \& thus did I the Lord God appoint unto man the days of his probation that by his natural death he might be raised in immortality unto eternal life even as many as would believe on my name
				\item[HSN] and thus did I the Lord God appoint unto man the days of his probation that by his natural death he might be raised in immortality unto eternal life even as many as would believe
				% \item[EMS] 
				% \item[BC] 
				% \item[]
			\end{compactdesc}
		\end{framed}

	% -----
	\subsubsection{Section 29:44} % *DC29-44 
	% -----
		\label{DC29-44}

		And they that believe not unto eternal damnation; for they cannot be redeemed from their spiritual fall, because they repent not;

		\begin{framed}
		\scriptsize
			\begin{compactdesc}
				\item[JSP] \& they that believe not unto eternal damnation for they cannot be redeemed from their spiritual fall Because they repent not
				\item[HSN] \& they that believe not unto eternal damnation for they cannot be redeemed from their spiritual fall because they repented not
				% \item[EMS] 
				% \item[BC] 
				% \item[]
			\end{compactdesc}
		\end{framed}

	% -----
	\subsubsection{Section 29:45} % *DC29-45 
	% -----
		\label{DC29-45}

		For they love darkness rather than light, and their deeds are evil, and they receive their wages of whom they \hyperref[LIST]{list} to obey.%
			\footnote{%
				% \label{}%
				Compare \hyperref[MOSIAH-2-33]{Mosiah 2:33} and \hyperref[ALMA-3-27]{Alma 3:27}.
				}

		\begin{framed}
		\scriptsize
			\begin{compactdesc}
				\item[JSP] for they love darkness more than light \& their deeds are evil \& they receive their wages of whom they list to obey
				\item[HSN] for they will Love darkness rather than light and their deeds are evil \& they receve their wages of whom they list to obey
				% \item[EMS] 
				% \item[BC] 
				% \item[]
			\end{compactdesc}
		\end{framed}

	% -----
	\subsubsection{Section 29:46} % *DC29-46 
	% -----
		\label{DC29-46}

		But behold, I say unto you, that little children are redeemed from the foundation of the world through mine Only Begotten;

		\begin{framed}
		\scriptsize
			\begin{compactdesc}
				\item[JSP] But Behold I say unto you that little children are redeemed from the foundation of the \sout{word} world through mine only begotten
				\item[HSN] but behold I say unto you that little child[re]n are redeemed from the \sout{fall} foundation of the world through mine onely begotten
				% \item[EMS] 
				% \item[BC] 
				% \item[]
			\end{compactdesc}
		\end{framed}

	% -----
	\subsubsection{Section 29:47} % *DC29-47 
	% -----
		\label{DC29-47}

		Wherefore, they cannot sin, for power is not given unto Satan to \hyperref[TEMPTATION]{tempt} little children, until they begin to become%
			\footnote{%
				% \label{}%
				Note that the JSP and HSN versions have ``be accountable,'' not ``become accountable.''
				}
		\hyperref[ACCOUNTABILITY]{accountable} before me;

		\begin{framed}
		\scriptsize
			\begin{compactdesc}
				\item[JSP] Wherefore they cannot sin for power is not given to Satan to tempt little children until they begin to be accountable before me
				\item[HSN] wherefore they cannot sin for power is not given unto Satan to tempt little children until they begin to be accountable befor me
				% \item[EMS] 
				% \item[BC] 
				% \item[]
			\end{compactdesc}
		\end{framed}

	% -----
	\subsubsection{Section 29:48} % *DC29-48 
	% -----
		\label{DC29-48}

		For it is given unto them even as I will, according to mine own pleasure, that great things may be required at the hand of their fathers.

		\begin{framed}
		\scriptsize
			\begin{compactdesc}
				\item[JSP] for it is given unto them even as I will according to mine own \sout{will} pleasure that great things may be required at the hand of their fathers
				\item[HSN] for it is given unto them even as I will according to mine own pleasure that great things may be required at the hand of their fathers
				% \item[EMS] 
				% \item[BC] 
				% \item[]
			\end{compactdesc}
		\end{framed}

	% -----
	\subsubsection{Section 29:49} % *DC29-49 
	% -----
		\label{DC29-49}

		And, again, I say unto you, that whoso having knowledge, have I not commanded to repent?

		\begin{framed}
		\scriptsize
			\begin{compactdesc}
				\item[JSP] \& again I say unto you that whoso having knowledge have not I commanded to Repent
				\item[HSN] \& again I say unto you that whoso having knowledge have I not commanded to repent
				% \item[EMS] 
				% \item[BC] 
				% \item[]
			\end{compactdesc}
		\end{framed}

	% -----
	\subsubsection{Section 29:50} % *DC29-50 
	% -----
		\label{DC29-50}

		And he that hath no understanding, it remaineth in me to do according as it is written.  And now I declare no more unto you at this time.  Amen.

		\begin{framed}
		\scriptsize
			\begin{compactdesc}
				\item[JSP] \& he that hath no understanding it remaineth in me to do according as it is written \& now behold I declare no more unto you at this time amen
				\item[HSN] \& he that hath no understanding it remaineth in me to do accord[in]g as it is written \& now I declare no more \sout{at} unto you at this time Amen
				% \item[EMS] 
				% \item[BC] 
				% \item[]
			\end{compactdesc}
		\end{framed}